%============================================================
% CH1 · PATTERN 1 — DEGREES AND COUNTING: THEORY
%============================================================
\subsection{Pattern Theory}
\label{ssec:ch1-p1-theory}

\begin{patternbox}{Pattern 1 at a Glance}
Many graph questions are \textbf{local}: the answer for a vertex depends only on the edges that touch it. Then we never walk along the graph --- we read the edge list \textbf{once} and keep a counter per vertex.

\begin{center}
\begin{tikzpicture}[x=1.15cm, y=1cm]
    \node[gv] (a) at (0,0) {0};
    \node[gv] (b) at (1.5,0.8) {1};
    \node[gv] (c) at (1.5,-0.8) {2};
    \node[gv] (d) at (3,0.8) {3};
    \node[gv] (e) at (4.5,0.8) {4};
    \node[gv] (f) at (4.5,-0.8) {5};
    \draw[ge] (a) -- (b); \draw[ge] (a) -- (c); \draw[ge] (b) -- (c);
    \draw[ge] (b) -- (d); \draw[ge] (d) -- (e);
    \node[gdist, left=1mm of a] {2};
    \node[gdist, above=1mm of b] {3};
    \node[gdist, below=1mm of c] {2};
    \node[gdist, above=1mm of d] {2};
    \node[gdist, above=1mm of e] {1};
    \node[gdist, right=1mm of f] {0};
    \node[note, anchor=west, text width=4.4cm, align=left] at (5.5,0)
        {green numbers = degrees, found in one pass over the 5 edges: $2+3+2+2+1+0 = 10 = 2\cdot 5$};
\end{tikzpicture}
\end{center}

\begin{recurrencebox}
\centering
for every edge $(u,v)$:\quad $\deg(u){+}{+},\ \deg(v){+}{+}$ \quad (directed: $\outdeg{u}{+}{+},\ \indeg{v}{+}{+}$)\\[3pt]
then one loop over the vertices answers the question\\[2pt] $\bigO{|V|+|E|}$ time, \ $\bigO{|V|}$ memory
\end{recurrencebox}
Typical answers: ``the vertex with degree $n-1$'', ``in-degree $n-1$ and out-degree $0$'', ``$\deg(u)+\deg(v)$ minus a correction'', ``$\sum_v w(v)\cdot\deg(v)$''. No adjacency list, no DFS, no BFS.
\end{patternbox}

%------------------------------------------------------------
\subsubsection*{From Part A to Pattern 1}

Lesson~\ref{ssec:ch1-vocabulary} defined the \textbf{degree} $\deg(v)$ of a vertex (the number of edges touching it) and, for directed graphs, the \textbf{in-degree} $\indeg{v}$ and \textbf{out-degree} $\outdeg{v}$. Lesson~\ref{ssec:ch1-storing} showed how to store a whole graph. This pattern is the observation that for a surprising number of problems \emph{you do not need to store the graph at all}: an array \texttt{deg[]} of $|V|$ integers is enough. It is the simplest graph algorithm there is, and it is the first thing to try before writing any traversal.

\begin{diagrambox}
\begin{tikzpicture}[x=1cm, y=1cm]
    \node[font=\small\bfseries, text=deepblue] at (0,2.25) {edge list};
    \foreach \u/\v [count=\i] in {0/1, 0/2, 1/2, 1/3, 3/4}{
        \node[cell, minimum width=13mm] (E\i) at (0,2.1-0.5*\i) {(\u, \v)};
    }
    \draw[flow] (1,0.8) -- node[lbl, above]{one pass} (3.2,0.8);
    \node[font=\small\bfseries, text=deepblue] at (6.1,2.25) {\texttt{deg[v]}};
    \foreach \d [count=\i from 0] in {2,3,2,2,1,0}{
        \node[idx] at (3.85+\i*0.9,1.65) {\i};
        \node[compcell] (D\i) at (3.85+\i*0.9,1.05) {\d};
    }
    \node[note, text width=5.6cm] at (6.1,0.05) {every edge adds 1 to \emph{each} of its two ends --- the edge $(1,3)$ is one of the three units in \texttt{deg[1]} and one of the two in \texttt{deg[3]}};
\end{tikzpicture}
\figcap{The whole ``algorithm'': an edge list goes in, a degree array comes out.}
\end{diagrambox}

%------------------------------------------------------------
\subsubsection*{What a degree can tell you}

\begin{table}[H]
\centering\small
\renewcommand{\arraystretch}{1.2}
\begin{tabularx}{\linewidth}{@{}>{\raggedright\arraybackslash}p{3.7cm} >{\raggedright\arraybackslash}p{3.5cm} X@{}}
\toprule
\textbf{Question about vertex $v$} & \textbf{Degree test} & \textbf{Typical wording} \\
\midrule
isolated vertex & $\deg(v) = 0$ & ``a city with no roads'' \\
leaf (end of a path, tree leaf) & $\deg(v) = 1$ & ``dead end'', ``leaf'' \\
connected to everybody & $\deg(v) = n-1$ (simple graph) & ``center of a star'', ``universal vertex'' \\
source of a directed graph & $\indeg{v} = 0$ & ``nobody points to it'', ``no prerequisites'' \\
sink of a directed graph & $\outdeg{v} = 0$ & ``trusts nobody'', ``terminal state'' \\
pointed to by all, points to none & $\indeg{v}=n-1,\ \outdeg{v}=0$ & ``judge'', ``celebrity'' (Problem~1.1.1) \\
every vertex looks the same & all $\deg(v)$ equal ($k$-regular) & ``everyone has exactly $k$ friends'' \\
\bottomrule
\end{tabularx}
\end{table}

For directed graphs we keep \emph{two} counters per vertex. An arrow $u \to v$ leaves $u$ and enters $v$:

\begin{diagrambox}
\begin{tikzpicture}[x=1.25cm, y=1cm]
    \node[gvsrc] (a) at (0,0.9) {0};
    \node[gv] (b) at (1.8,1.6) {1};
    \node[gv] (c) at (1.8,0) {2};
    \node[gvsrc] (d) at (0,-0.7) {3};
    \node[gvdst] (e) at (3.6,0) {4};
    \draw[gde] (a) -- (b); \draw[gde] (a) -- (c); \draw[gde] (b) -- (c);
    \draw[gde] (d) -- (c); \draw[gde] (c) -- (e);
    \foreach \n/\i/\o/\p in {a/0/2/left, b/1/1/above, c/3/1/below, d/0/1/left, e/1/0/below}{
        \node[gdist, \p=1.5mm of \n] {in~\i\ / out~\o};
    }
    \node[note, anchor=west, text width=5.6cm, align=left] at (4.9,0.5)
        {\textcolor{successgreen}{\textbf{green}} vertices 0 and 3: in-degree 0 (\emph{sources})\\
         \textcolor{accentorange}{\textbf{orange}} vertex 4: out-degree 0 (a \emph{sink})\\[3pt]
         $\sum \indeg{v} = \sum \outdeg{v} = |E| = 5$:\\ every arrow has exactly one tail and one head.};
\end{tikzpicture}
\figcap{In-degree and out-degree of every vertex, again from one pass over the edges.}
\end{diagrambox}

%------------------------------------------------------------
\subsubsection*{The one law of degrees: the handshake lemma}

Cut every edge in the middle. Each edge becomes two \emph{half-edges} (``stubs''), one attached to each end. A vertex of degree $d$ holds exactly $d$ stubs, so counting stubs vertex by vertex or edge by edge gives
\[
\sum_{v \in V} \deg(v) \;=\; 2\,|E| .
\]

\begin{diagrambox}
\begin{tikzpicture}[x=1.15cm, y=1cm]
    \node[gv] (a) at (0,0) {0};
    \node[gv] (b) at (1.5,0.8) {1};
    \node[gv] (c) at (1.5,-0.8) {2};
    \node[gv] (d) at (3,0.8) {3};
    \node[gv] (e) at (4.5,0.8) {4};
    \foreach \x/\y in {a/b, a/c, b/c, b/d, d/e}{
        \draw[draw=primaryblue, line width=1.6pt] (\x) -- ($(\x)!0.5!(\y)$);
        \draw[draw=accentorange, line width=1.6pt] ($(\x)!0.5!(\y)$) -- (\y);
        \fill[white] ($(\x)!0.5!(\y)$) circle (1.8pt);
    }
    \node[note, anchor=west, text width=6cm, align=left] at (5.3,0)
        {5 edges $\Rightarrow$ 10 stubs (blue + orange).\\ Vertex 1 holds 3 of them, vertex 4 holds 1, \ldots\\[3pt]
        \textbf{Consequences:} the degree sum is even; the number of odd-degree vertices is even; $n$ vertices of degree $\ge k$ force at least $nk/2$ edges.};
\end{tikzpicture}
\figcap{Counting the same stubs twice proves the handshake lemma.}
\end{diagrambox}

The lemma is a free consistency check (``can these numbers be the degrees of a graph?'') and the reason why the counting loop runs in $\bigO{|V|+|E|}$: it does exactly $2|E|$ increments. For a directed graph the analogue is $\sum_v \indeg{v} = \sum_v \outdeg{v} = |E|$.

%------------------------------------------------------------
\subsubsection*{When to recognise it --- and when not}

Use Pattern 1 when
\begin{itemize}[leftmargin=*]
    \item the answer for a vertex depends only on \textbf{the edges touching it} (``trusts nobody'', ``connected to all'', ``number of roads at a city'');
    \item the answer for a \textbf{pair} of vertices depends only on their two degrees and on whether they are adjacent (Problem~1.1.2);
    \item the quantity is a \textbf{sum over edges} of something that depends on the endpoints: regroup it as a sum over vertices weighted by degree (Exercise~E1.3);
    \item the input \textbf{promises} a shape (``the graph is a star'', ``at most one person is trusted by all'') and you only need to locate the special vertex.
\end{itemize}

\textbf{Do NOT use it} for anything that depends on \emph{paths}: reachability, connectivity, distances, cycles. Degrees are blind to how the edges are wired together:

\begin{diagrambox}
\begin{tikzpicture}[x=1cm, y=1cm]
    \begin{scope}
        \foreach \i in {0,...,5}{ \node[gvA] (h\i) at ({90-60*\i}:1.1) {\i}; }
        \foreach \i [evaluate=\i as \j using {int(mod(\i+1,6))}] in {0,...,5}{ \draw[ge] (h\i) -- (h\j); }
        \node[font=\small\bfseries] at (0,-1.75) {one cycle of length 6: connected};
    \end{scope}
    \begin{scope}[xshift=6.2cm]
        \node[gvB] (t0) at (-1.5,0.7) {0}; \node[gvB] (t1) at (-0.75,-0.5) {1}; \node[gvB] (t2) at (-2.25,-0.5) {2};
        \node[gvB] (s0) at (1.5,0.7) {3}; \node[gvB] (s1) at (2.25,-0.5) {4}; \node[gvB] (s2) at (0.75,-0.5) {5};
        \draw[ge] (t0) -- (t1) -- (t2) -- (t0);
        \draw[ge] (s0) -- (s1) -- (s2) -- (s0);
        \node[font=\small\bfseries] at (0,-1.75) {two triangles: disconnected};
    \end{scope}
    \node[note] at (3.1,1.55) {both graphs: 6 vertices, 6 edges, every degree $= 2$};
\end{tikzpicture}
\figcap{Identical degree sequences, different answers to ``is the graph connected?''. Such questions need a traversal (Pattern~2).}
\end{diagrambox}

%------------------------------------------------------------
\subsubsection*{The algorithm and why it is correct}

\begin{enumerate}[leftmargin=*]
    \item \step{Allocate} \texttt{deg[0..n]} (one extra slot so that 1-indexed labels fit), all zero. Directed: \texttt{in[]} and \texttt{out[]}, or a single \texttt{score[] = in - out}.
    \item \step{Scan} the edge list once; for $(u,v)$ do \texttt{deg[u]++, deg[v]++} (directed: \texttt{out[u]++, in[v]++}).
    \item \step{Answer} with a loop over the vertices (or over the pairs, or over the edges again).
\end{enumerate}

\begin{invariantbox}[Invariant of the counting loop]
After the first $i$ edges $e_1,\dots,e_i$ have been processed, \texttt{deg[v]} equals the number of edges among $e_1,\dots,e_i$ that have $v$ as an endpoint (a self-loop $(v,v)$ counts twice).

\textbf{Proof.} True for $i = 0$ (all zero, no edges). Processing $e_{i+1} = (u,w)$ adds 1 to exactly the counters of its endpoints $u$ and $w$ --- the only vertices whose count of incident edges grows --- and leaves all other counters unchanged. After the last edge, \texttt{deg[v]} $=\deg(v)$. \hfill$\blacksquare$

\textbf{Cost.} $2|E|$ increments plus $\bigO{|V|}$ to initialise and scan: $\bigO{|V|+|E|}$ time, $\bigO{|V|}$ extra memory --- optimal, since we must at least read the input.
\end{invariantbox}

%------------------------------------------------------------
\subsubsection*{Four counting tools}

\textbf{Tool 1 --- one counter per vertex.} The table above. Find the vertices whose counter has a special value.

\textbf{Tool 2 --- a net score.} When the condition is ``in-degree exactly $A$ \emph{and} out-degree exactly $0$'', one array \texttt{score[v] = in[v] - out[v]} often suffices: if in-degrees can never exceed $A$, then $\texttt{score}[v] = A \iff \indeg{v} = A$ and $\outdeg{v} = 0$. This is the whole solution of Problem~1.1.1.

\textbf{Tool 3 --- pairs: inclusion--exclusion.} The number of edges touching $u$ \emph{or} $v$ is
\[
|\,E(u) \cup E(v)\,| \;=\; \deg(u) + \deg(v) - [\,u \text{ and } v \text{ are adjacent}\,],
\]
because the edge $(u,v)$, if present, was counted once at $u$ and once at $v$. Degrees plus an adjacency test answer every such pair question (Problem~1.1.2).

\begin{diagrambox}
\begin{tikzpicture}[x=1.2cm, y=1cm]
    \node[gvsel] (u) at (0,0) {$u$};
    \node[gvsel] (v) at (2.4,0) {$v$};
    \node[gv] (a) at (-1.3,0.8) {}; \node[gv] (b) at (-1.3,-0.8) {};
    \node[gv] (c) at (3.7,0.8) {}; \node[gv] (d) at (3.7,-0.8) {}; \node[gv] (e) at (3.7,0) {};
    \draw[draw=primaryblue, line width=1.3pt] (u) -- (a); \draw[draw=primaryblue, line width=1.3pt] (u) -- (b);
    \draw[draw=accentorange, line width=1.3pt] (v) -- (c); \draw[draw=accentorange, line width=1.3pt] (v) -- (d);
    \draw[draw=accentorange, line width=1.3pt] (v) -- (e);
    \draw[gpath, draw=softpurple] (u) -- node[gw, above=1mm]{counted twice} (v);
    \node[note, anchor=west, text width=4.6cm, align=left] at (4.3,0)
        {$\deg(u) = 3,\ \deg(v) = 4$\\ edges touching $u$ or $v$: $3 + 4 - 1 = 6$\\ (without the purple edge: $2 + 3 - 0 = 5$)};
\end{tikzpicture}
\figcap{Inclusion--exclusion for two vertices: the shared edge is subtracted once.}
\end{diagrambox}

\textbf{Tool 4 --- turn a sum over edges into a sum over vertices.} If every edge contributes $f(u) + f(v)$, then vertex $v$ appears once per incident edge, so
\[
\sum_{(u,v) \in E} \big(f(u) + f(v)\big) \;=\; \sum_{v \in V} f(v)\cdot \deg(v).
\]
In the same spirit the number of \emph{paths with two edges} $u - v - w$ (``cherries'') is $\sum_v \binom{\deg(v)}{2}$: choose the middle vertex, then two of its edges.

\begin{diagrambox}
\begin{tikzpicture}[x=1.15cm, y=1cm]
    \node[gv] (a) at (0,0) {0};
    \node[gvcur] (b) at (1.5,0.8) {1};
    \node[gv] (c) at (1.5,-0.8) {2};
    \node[gv] (d) at (3,0.8) {3};
    \node[gv] (e) at (4.5,0.8) {4};
    \draw[gpath] (a) -- (b); \draw[gpath] (b) -- (c); \draw[gpath] (b) -- (d);
    \draw[ge] (a) -- (c); \draw[ge] (d) -- (e);
    \node[note, anchor=west, text width=6.2cm, align=left] at (5.3,0)
        {the 3 orange edges at vertex 1 give $\binom{3}{2} = 3$ cherries with middle 1: $0{-}1{-}2,\ 0{-}1{-}3,\ 2{-}1{-}3$.\\[2pt]
        All cherries: $\binom22+\binom32+\binom22+\binom22+\binom12 = 1+3+1+1+0 = 6$.};
\end{tikzpicture}
\figcap{Counting small structures by their centre: a cherry is fixed by its middle vertex and two of its edges.}
\end{diagrambox}

%------------------------------------------------------------
\subsubsection*{Variations map}

\begin{diagrambox}
\begin{tikzpicture}[
    var/.style={rectangle, rounded corners=2pt, draw=patternborder, fill=patternbg,
                font=\scriptsize, align=center, text width=3.25cm, minimum height=1.25cm, inner sep=3pt}]
    \node[var, fill=answerfill, draw=accentorange, text width=5cm, font=\scriptsize\bfseries] (root) at (0,0)
        {Degree counting\\one pass: \texttt{deg[u]++, deg[v]++}};
    \node[var] (v1) at (-5.4,-2) {\textbf{special vertex}\\star center, universal vertex, leaf, isolated (E1.1)};
    \node[var] (v2) at (-1.8,-2) {\textbf{in/out-degree}\\judge, sources, sinks (Problem~1.1.1)};
    \node[var] (v3) at (1.8,-2)  {\textbf{pairs}\\$\deg u + \deg v - [uv \in E]$ (Problem~1.1.2)};
    \node[var] (v4) at (5.4,-2)  {\textbf{sum over edges}\\$= \sum_v f(v)\deg(v)$; sort degrees (E1.3)};
    \node[var] (v5) at (-3.6,-3.6) {\textbf{best neighbours}\\top-$k$ neighbour values per vertex (E1.4)};
    \node[var] (v6) at (0,-3.6)    {\textbf{parity \& existence}\\handshake: is it a degree sequence? (E1.5)};
    \node[var] (v7) at (3.6,-3.6)  {\textbf{build from degrees}\\a tree with given degrees (E1.7)};
    \node[draw=patternborder, dashed, rounded corners=4pt, fit=(v1)(v4)(v5)(v7), inner sep=5pt] (grp) {};
    \draw[flow] (root.south) -- (grp.north);
\end{tikzpicture}
\figcap{Every box is solved by degrees alone. Degree tricks on harder problems (triangles, ordering by degree) return in Chapter~22.}
\end{diagrambox}

\textbf{Twists and how to handle them}
\begin{itemize}[leftmargin=*]
    \item \textbf{Labels $1..n$:} allocate $n+1$ slots and ignore slot 0, instead of shifting every label.
    \item \textbf{Self-loops} $(v,v)$ add 2 to $\deg(v)$ (two stubs); \textbf{multi-edges} add 1 each time. If the problem says ``number of \emph{neighbours}'', deduplicate first (a set per vertex, or sort the edge list).
    \item \textbf{Duplicates in a directed list} break the net-score trick: $\indeg{v}$ could exceed $n-1$. Read the constraints (``all pairs are unique'').
    \item \textbf{Large sums:} $\sum_v f(v)\deg(v)$ easily exceeds $2^{31}$ --- use \texttt{long long}.
    \item \textbf{Pairs with $n$ up to $10^5$:} the $\bigO{n^2}$ pair loop is too slow; only vertices of the largest one or two degrees can win (Approach 3 of Problem~1.1.2).
\end{itemize}

%------------------------------------------------------------
\subsubsection*{Code templates}

\begin{lstlisting}[style=cpp]
#include <bits/stdc++.h>
using namespace std;
typedef long long ll;

// Input: n m, then m edges "u v" (1-indexed). Prints a few degree facts.
int main() {
    int n, m;
    if (!(cin >> n >> m)) return 0;
    vector<int> deg(n + 1, 0);                  // undirected view
    vector<int> in(n + 1, 0), out(n + 1, 0);    // directed view (u -> v)
    for (int i = 0; i < m; i++) {
        int u, v;
        cin >> u >> v;
        deg[u]++; deg[v]++;                     // (1) both ends
        out[u]++; in[v]++;                      // (2) tail and head
    }
    // (3) special vertices
    for (int v = 1; v <= n; v++) {
        if (deg[v] == 0) cout << v << " is isolated\n";
        if (deg[v] == 1) cout << v << " is a leaf\n";
        if (deg[v] == n - 1) cout << v << " touches every other vertex\n";
        if (in[v] == 0) cout << v << " is a source\n";
        if (out[v] == 0) cout << v << " is a sink\n";
    }
    // (4) handshake check and cherries (paths with two edges)
    ll sum = 0, cherries = 0;
    for (int v = 1; v <= n; v++) {
        sum += deg[v];
        cherries += (ll)deg[v] * (deg[v] - 1) / 2;
    }
    cout << "degree sum " << sum << " = 2*" << m << ", cherries " << cherries << "\n";
}
\end{lstlisting}

The second template answers \emph{pair} questions. For $n$ up to a few thousand an adjacency matrix gives an $\bigO{1}$ adjacency test; for large $n$ use a hash set of normalised pairs $(\min, \max)$, as here.

\begin{lstlisting}[style=cpp]
#include <bits/stdc++.h>
using namespace std;
typedef long long ll;

struct PairCounter {
    int n;
    vector<int> deg;
    unordered_set<ll> edgeSet;                   // key = min * n + max (0-indexed)
    PairCounter(int n_) : n(n_), deg(n_, 0) {}
    ll key(int u, int v) const { if (u > v) swap(u, v); return (ll)u * n + v; }
    void addEdge(int u, int v) {
        deg[u]++; deg[v]++;
        edgeSet.insert(key(u, v));
    }
    bool adjacent(int u, int v) const { return edgeSet.count(key(u, v)) > 0; }
    // number of edges touching u or v (Tool 3)
    int touching(int u, int v) const { return deg[u] + deg[v] - (adjacent(u, v) ? 1 : 0); }
};

int main() {
    PairCounter g(6);
    int E[5][2] = {{0, 1}, {0, 2}, {1, 2}, {1, 3}, {3, 4}};
    for (auto& e : E) g.addEdge(e[0], e[1]);
    cout << g.touching(0, 1) << " " << g.touching(1, 4) << " " << g.touching(4, 5) << "\n";  // 4 4 1
}
\end{lstlisting}

\begin{insightbox}[Check the templates by hand]
On the running example ($6$ vertices, edges $01, 02, 12, 13, 34$; degrees $2,3,2,2,1,0$): \texttt{touching(0,1)} $= 2+3-1 = 4$ (edges $01,02,12,13$); \texttt{touching(1,4)} $= 3+1-0 = 4$ (edges $01,12,13,34$); \texttt{touching(4,5)} $= 1+0-0 = 1$. Fed with \texttt{6 5} and the same edges shifted to labels $1..6$, the first template reports vertex 6 as isolated, vertex 5 as a leaf, vertices 1 and 6 as sources, vertices 3, 5, 6 as sinks, the degree sum $10 = 2\cdot5$ and $6$ cherries.
\end{insightbox}

%------------------------------------------------------------
\subsubsection*{Common mistakes}

\begin{warningbox}[Mistake 1 --- counting only one endpoint]
\begin{lstlisting}[style=snippet]
for (auto& e : edges) deg[e[0]]++;                  // WRONG for undirected graphs
for (auto& e : edges) { deg[e[0]]++; deg[e[1]]++; } // RIGHT
\end{lstlisting}
Sanity check: the sum of all degrees must be $2|E|$.
\end{warningbox}

\begin{warningbox}[Mistake 2 --- off-by-one in the labels]
Labels $1..n$ with \texttt{vector<int> deg(n)} write past the end at label $n$ (undefined behaviour, often a silent wrong answer). Allocate $n+1$, or subtract 1 when reading --- consistently.
\end{warningbox}

\begin{warningbox}[Mistake 3 --- using degrees for a path question]
``Every vertex has degree $\ge 1$, so the graph is connected'' is false (two separate edges). ``$|E| = n-1$, so it is a tree'' is false too (a triangle plus an isolated vertex has $n=4$, $|E|=3$). Connectivity needs a traversal (Pattern~2) or DSU (Chapter~6).
\end{warningbox}

\begin{warningbox}[Mistake 4 --- double-counting the shared edge]
For a pair $(u,v)$, $\deg(u)+\deg(v)$ counts the edge $(u,v)$ twice. Subtract it --- but only if it exists.
\end{warningbox}

\begin{warningbox}[Mistake 5 --- forgetting the empty and the tiny cases]
With $n = 1$ and no edges, the only vertex has degree $0 = n-1$: it \emph{is} connected to all (zero) others. With no edges at all, every counter is 0 --- make sure the code still returns a valid answer instead of reading \texttt{edges[0]}.
\end{warningbox}

\begin{checklistbox}[``I'm stuck'' checklist for Pattern 1]
\begin{itemize}[leftmargin=*]
    \item[$\square$] Is the question \emph{local}? Can I state the condition using only the edges at one vertex (or at two)?
    \item[$\square$] Directed or undirected? Do I need in-degree, out-degree, both, or their difference?
    \item[$\square$] Is it a sum over edges? Rewrite it as $\sum_v f(v)\deg(v)$.
    \item[$\square$] For pairs: did I subtract the shared edge? Is $\bigO{n^2}$ over pairs affordable, or must I restrict to the top degrees?
    \item[$\square$] Self-loops, multi-edges, duplicates: what do the constraints promise?
    \item[$\square$] Does the answer depend on paths after all? Then go to Pattern~2.
\end{itemize}
\end{checklistbox}
